% =============================================================================
% CONNECTIONS TO EXISTING FRAMEWORKS (1-2 pages)
%
% Spec requirements:
%   - Empowerment maximization (formal bridge)
%   - Sen's capability approach (philosophical bridge)
%   - Free energy principle (structural bridge)
%   - Constitutional AI / RLHF (practical bridge)
%
% Reference material:
%   - Klyubin et al. (2005), Sen (1999), Friston (2013)
%   - Christiano et al. (2017), Bai et al. (2022)
% =============================================================================

\section{Connections to Existing Frameworks}
\label{sec:connections}

\subsection{Empowerment Maximization}

% TODO: Empowerment = MI between actions and future states (Klyubin 2005)
% GFM is social empowerment: across joint action space of all agents
% Key difference: empowerment is agent-centric, GFM is population-centric

\subsection{Sen's Capability Approach}

% TODO: Well-being measured by available functionings, not happiness/resources
% GFM: alignment measured by achievable goals, not fixed utility
% Key difference: capability approach is descriptive, GFM is prescriptive

\subsection{Free Energy Principle}

% TODO: FEP (Friston 2013): minimize prediction error
% FREE: maximize goal frontier of society
% Structural parallel: coherent world models + adaptation
% "Free Relative Energy Ethics" is a deliberate echo

\subsection{Constitutional AI and RLHF}

% TODO: Human feedback as noisy channel for communicating goal frontier
% GFM predicts where RLHF fails: narrow evaluator pool, reward hacking
